\chapter{Conclusiones y trabajo futuro}
\label{cap:conclusiones}

Esta memoria preguntó si arquitecturas tipo \Transformer\ pueden operar directamente sobre series de tiempo no equiespaciadas y si el tiempo físico aporta información identificable. La primera pregunta queda respondida a nivel de implementación y factibilidad: el pipeline conserva eventos irregulares, formula consultas en timestamps arbitrarios y entrena sin remuestreo a una grilla fija. La segunda exige más cautela. El benchmark legado no separa completamente horizonte físico, orden de consulta y densidad de eventos; la evaluación física congelada fue diseñada para hacerlo, pero sólo su campaña completa puede establecer el ranking correspondiente.

\section{Conclusiones principales}
\label{sec:conclusiones-principales}

Las conclusiones se organizan por nivel de evidencia:
\begin{enumerate}
    \item \textbf{Factibilidad técnica.} \CTTransformer\ y sus variantes procesan historias irregulares, múltiples targets y datos multivariados asincrónicos sin imputar una malla temporal completa. Los timestamps absolutos permanecen en \textttbreak{float64} hasta restar un origen local, evitando que gaps pequeños colapsen por un cast prematuro.
    \item \textbf{Resultado legado descriptivo.} En cinco series reales, el encoder--decoder pequeño con ajuste AR contiguo obtiene las menores medias publicadas: MSE 0.1979, RMSE 0.4417 y MAE 0.2193, con victoria en 5/5 datasets. Las semillas se promedian dentro de cada dataset; con cinco unidades no corresponde declarar significancia estadística ni universalidad.
    \item \textbf{Límite de la ablación posicional.} Que el modelo sin encoding continuo casi empate al encoder base no demuestra que el tiempo sea irrelevante. Los offsets están medidos en eventos, las consultas conservan orden y la historia contiene proxies de densidad. La tarea puede resolverse parcialmente mediante esas señales, persistencia o dinámica ordinal.
    \item \textbf{Arquitecturas para tiempo físico.} \QueryCross\ codifica sólo la historia y resuelve queries independientes mediante atención cruzada con lag relativo, identidad de sensor y features temporales. \ContinuousBasis\ representa el pronóstico como una función continua del horizonte. Ambas eliminan el embedding de slot objetivo y admiten \CTSSM.
    \item \textbf{Dinámica continua eficiente.} \CTSSM\ aplica transiciones diagonales estables dependientes del gap y una actualización gated mediante un scan paralelo asociativo. Es una alternativa liviana, sin solver ODE, pero no representa toda la clase de modelos de estado continuo.
    \item \textbf{Extensión probabilística.} Las cabezas gaussianas estiman media y escala heteroscedástica, se seleccionan por NLL de validación y se evalúan con NLL, CRPS y cobertura. Una buena cobertura en estos procesos no garantiza calibración fuera del dominio sintético ni captura futuros multimodales.
    \item \textbf{Protocolo físico auditable.} La campaña final fija historia y horizonte en unidades físicas, obtiene targets desde una trayectoria limpia, aleatoriza queries, incluye controles de identificabilidad y reporta estratos de dificultad. Los hiperparámetros se congelan en YAML y test queda fuera de selección.
\end{enumerate}

\section{Contraste con la hipótesis}
\label{sec:contraste-hipotesis}

La hipótesis de factibilidad queda respaldada: es posible entrenar Transformers sobre eventos irregulares y consultar tiempos futuros explícitos. El benchmark legado aporta evidencia descriptiva de competitividad y muestra la utilidad del identificador historia--objetivo, pero no identifica el efecto aislado del tiempo físico.

La hipótesis más fuerte---que modelar lags físicos mejora de manera reproducible frente a controles sin tiempo---queda condicionada a los cuatro reportes generados por la evaluación física completa. Los smoke tests realizados durante desarrollo comprueban contratos de datos, forward, backward y escritura de artefactos; no tienen cobertura ni presupuesto para seleccionar una arquitectura. Si las tablas condicionales del Capítulo~\ref{cap:resultados} no están presentes, esta memoria no afirma un ganador entre \QueryCross, \ContinuousBasis\ y sus variantes.

La evaluación física es congelada y versionada, no una prueba ciega: sus procesos sintéticos y diagnósticos fueron inspeccionados mientras se desarrollaba el sistema. Sus resultados describen los 16 datasets principales y las dos unidades de estrés predefinidas; no establecen validez externa en series reales.

\section{Limitaciones}
\label{sec:limitaciones-finales}

Las conclusiones deben leerse bajo las siguientes limitaciones:
\begin{itemize}
    \item El benchmark legado contiene sólo cinco datasets independientes. Tratar tres semillas por dataset como $n=15$ produciría pseudorreplicación; por eso se retiraron sus intervalos y pruebas a nivel de corrida.
    \item El legado define historias y horizontes por número de eventos, mientras el protocolo físico los define por duración. Sus métricas tienen estimandos distintos y no son comparables numéricamente.
    \item La evaluación física usa procesos sintéticos. Permite targets contrafactuales limpios y corruptelas controladas, pero no representa toda la diversidad, missingness informativo ni cambios de régimen de datos reales.
    \item Los presets físicos fueron observados durante el desarrollo. Congelar configuración y stress set reduce grados de libertad posteriores, pero no elimina el riesgo de adaptación al benchmark.
    \item La familia gaussiana supone una distribución unimodal condicional. NLL y cobertura pueden verse afectadas por especificación distributiva, tamaño muestral y escalado.
    \item Los hiperparámetros físicos se fijan a priori para comparabilidad y costo; no se demuestra optimalidad global de ninguna familia.
\end{itemize}

\section{Trabajo futuro}
\label{sec:trabajo-futuro}

Las continuaciones prioritarias son:
\begin{enumerate}
    \item \textbf{Completar y archivar la campaña física}: ejecutar \textttbreak{run} y \textttbreak{report} con las tres semillas y verificar que los cuatro includes LaTeX provengan de las 16 unidades principales, manteniendo stress separado.
    \item \textbf{Validación externa preregistrada}: fijar hipótesis, métricas y presupuesto antes de evaluar nuevas colecciones reales no usadas durante desarrollo.
    \item \textbf{Aumentar unidades independientes}: incorporar más procesos y datasets, preservando la agregación semillas $\rightarrow$ dataset y aplicando corrección por comparaciones múltiples cuando corresponda.
    \item \textbf{Estudiar missingness informativo}: añadir controles en los que la ausencia de una medición dependa del estado y comparar máscaras, tokens de ausencia y modelos generativos conjuntos.
    \item \textbf{Ampliar incertidumbre}: contrastar la cabeza gaussiana con cuantiles, mezclas o distribuciones robustas, evaluando calibración condicional por horizonte y régimen de gaps.
    \item \textbf{Perfilar eficiencia}: medir memoria y latencia de \QueryCross, \ContinuousBasis\ y \CTSSM\ al variar longitud de historia, número de sensores y cantidad de consultas.
\end{enumerate}

\section{Cierre}
\label{sec:cierre-memoria}

La contribución central no es afirmar que un encoding sinusoidal continuo resuelva por sí solo el muestreo irregular. El resultado metodológico es mostrar por qué esa afirmación no era identificable en el benchmark legado y construir una alternativa auditable: consultas independientes en tiempo físico, kernels de lag, bases continuas, estado estable, precisión temporal explícita y controles emparejados.

En consecuencia, esta memoria entrega una arquitectura extensible y un protocolo reproducible para estudiar el tiempo irregular con mayor rigor. El benchmark legado demuestra factibilidad y ofrece un resultado descriptivo sobre cinco series reales; la campaña física, cuando esté completa, permitirá comparar las nuevas variantes dentro de los procesos sintéticos cubiertos, sin confundir un smoke test con evidencia de ranking ni extrapolar más allá de ese dominio.
